\chapter{Problem Statement and Objectives}
\label{ch:problem_objectives}

\section{Problem Statement}
Practical ZK-Rollup performance is dictated by several interacting and often competing factors, including batch formation limits, sequencer scheduling algorithms, data availability footprint, Layer-1 submission costs, and proof-generation overhead \cite{habib2025bottlenecks,chaliasos2024benchmarking}. These factors do not affect performance independently; for instance, increasing the batch size reduces the amortized Layer-1 cost per transaction but simultaneously increases the queueing delay for early transactions in the batch.

The core research problem addressed in this report is the lack of a compact, modular, and reproducible experimental framework designed specifically to isolate how these batching, sequencing, proof-delay, and data availability choices affect ZK-Rollup throughput, latency, finalization time, and cost behavior under strictly comparable workloads.

\section{Research Questions}
This work is guided by the following primary research questions:
\begin{enumerate}
    \item How can a modular ZK-Rollup prototype be designed to support the controlled benchmarking of batching, sequencing, state execution, proof handling, and data availability choices?
    \item How do variations in batching strategy, sequencing policy, and data availability mode (calldata vs. blobs) shift the end-to-end performance across throughput, latency, data footprint, and cost dimensions?
    \item How can parameterized proof-delay measurements help isolate proving-stage effects from sequencing, execution, and Layer-1 submission overhead?
\end{enumerate}

\section{Objectives}
The primary objective of this project is to design, implement, and systematically evaluate a modular ZK-Rollup research prototype using controlled synthetic workloads. 

The expected outcomes include:
\begin{itemize}
    \item A reproducible, containerized experimental framework.
    \item Empirical insights into the inherent throughput--latency--cost trade-offs present in the Layer-2 pipeline.
    \item An open-source benchmarking harness that provides a foundation for future scalability research.
\end{itemize}
